\documentclass{article}

\usepackage{amsmath,amssymb}
\usepackage{xurl}
\usepackage[hidelinks]{hyperref}
\setlength{\emergencystretch}{2em}

% Shared commands for notes in docs/paper-gaps/.

\DeclareUnicodeCharacter{2080}{\ifmmode{}_{0}\else\textsubscript{0}\fi}
\DeclareUnicodeCharacter{2081}{\ifmmode{}_{1}\else\textsubscript{1}\fi}
\DeclareUnicodeCharacter{2082}{\ifmmode{}_{2}\else\textsubscript{2}\fi}
\DeclareUnicodeCharacter{2099}{\ifmmode{}_{n}\else\textsubscript{n}\fi}

\newcommand{\Lean}{\textsc{Lean}}
\ExplSyntaxOn
\NewDocumentCommand{\leanid}{m}{
  \ifmmode
    \textsf{\detokenize{#1}}
  \else
    \group_begin:
    \urlstyle{sf}
    \tl_set:Nn \l_tmpa_tl {#1}
    \tl_replace_all:Nnn \l_tmpa_tl {\_} {_}
    \exp_args:NV \path \l_tmpa_tl
    \group_end:
  \fi
}
\ExplSyntaxOff
\providecommand{\tr}{\operatorname{tr}}
\newcommand{\tnleanIssueBase}{https://github.com/LionSR/TNLean/issues/}
\newcommand{\tnleanPrBase}{https://github.com/LionSR/TNLean/pull/}
\newcommand{\tnleanDocBase}{https://sirui-lu.com/TNLean/docs/}
\newcommand{\ghissue}[1]{\href{\tnleanIssueBase#1}{GitHub issue \##1}}
\newcommand{\ghpr}[1]{\href{\tnleanPrBase#1}{PR \##1}}
\newcommand{\doclink}[1]{\href{\tnleanDocBase#1.html}{\path{#1}}}

% Dirac notation and the identity operator, as in blueprint/src/macros/common.tex.
% \providecommand keeps note-local definitions authoritative.
\usepackage{dsfont}
\providecommand{\ket}[1]{|#1\rangle}
\providecommand{\bra}[1]{\langle #1|}
\providecommand{\braket}[2]{\langle #1 | #2 \rangle}
\providecommand{\ketbra}[2]{|#1\rangle\!\langle #2|}
\providecommand{\Id}{\mathds{1}}
\providecommand{\one}{\mathds{1}}
\providecommand{\C}{\mathbb{C}}
\providecommand{\MN}[1]{M_{#1}(\C)}
\providecommand{\MD}{M_D(\C)}

% Verdict marker: \gapnote{<kind>}{<status>}. Kind is the policy
% classification (clarification, local-correction, scope-restriction,
% unfaithful, false-source, open-gap); status is open, wip (elimination
% underway), resolved, or historical. Machine-read by the site index;
% prints a verdict line.
\newcommand{\gapnote}[2]{\par\noindent\textbf{Verdict:} #1 (#2).\par}


% Fallbacks keep this note compilable when a different command.tex is found first.
\providecommand{\leanid}[1]{\nolinkurl{#1}}
\providecommand{\gapnote}[2]{\par\noindent\textbf{Verdict:} #1 (#2).\par}

\title{The Deficiency of a Level in Hastings' Hierarchical Projection}
\author{}
\date{2026-10-23}

\begin{document}
\maketitle
\gapnote{local-correction}{resolved}

\section{The source passage}

Section~II of Hastings' \emph{Solving gapped Hamiltonians
    locally}~\cite{Hastings2006Solving} (arXiv:cond-mat/0508554v3,
\texttt{proj3new.tex}) builds a state close to a ground state $\Psi_a$ from
local spectral projections. There $J$ bounds the norms of the local terms,
$\lVert H_i\rVert\le J$ (line 180), each approximate term $M_i$ acts within
distance $l_{\mathrm{proj}}$ of site $i$ (lines 284--307), and
$l_{\mathrm{fac}}$ is the factorization length, the error of the
approximation being of order $e^{-l_{\mathrm{proj}}/l_{\mathrm{fac}}}$
(lines 344--349). For a pair operator $M=M_{k,k+1}$ the ground state
is an approximate eigenvector,
\begin{align}\label{eq:hastings05_hierarchical_projection_ground_state_residual}
    \bigl\lVert(M-c)\Psi_a\bigr\rVert\le\epsilon,
    \qquad c=\langle\Psi_a,M\Psi_a\rangle,
\end{align}
by eq.~(low) (lines 633--639). The projection $P^1$ onto the eigenvectors of
$M$ with eigenvalue in $[c-x,c+x]$ then satisfies (lines 645--657)
\begin{align}\label{eq:hastings05_hierarchical_projection_projection_expectation}
    x\,\langle\Psi_a,(1-P^1)\Psi_a\rangle\le\epsilon,
    \qquad\langle\Psi_a,P^1\Psi_a\rangle\ge1-b_1,
    \qquad b_1=\epsilon/x .
\end{align}
At the next level the source sets
$\Psi^1_a=P^1_kP^1_{k+2}\Psi_a/\lVert P^1_kP^1_{k+2}\Psi_a\rVert$ and states
(line 689)
\begin{align}\label{eq:hastings05_hierarchical_projection_normalized_projection_distance}
    \lVert\Psi^1_a-\Psi_a\rVert\le2b_1 .
\end{align}
From this it bounds the residual of $\Psi^1_a$ by $\epsilon+2b_1J\,l_{\mathrm{proj}}$
(lines 690--695) and obtains the recursion
$b_m=(\epsilon+2b_{m-1})^{1-1/h}$ with $h=\log_2N$ levels (lines 696--704).

\section{The deficiency in norm}

An expectation bound does not give~(\ref{eq:hastings05_hierarchical_projection_normalized_projection_distance}). If
$\langle\Psi,(1-P)\Psi\rangle\le b$ for an orthogonal projection $P$ and a unit
vector $\Psi$, then only
\begin{align}\label{eq:hastings05_hierarchical_projection_sqrt_deficiency}
    \lVert(1-P)\Psi\rVert=\langle\Psi,(1-P)\Psi\rangle^{1/2}\le\sqrt b ,
\end{align}
and a normalized image can lie at distance of order $\sqrt b$, not $b$. Read
with~(\ref{eq:hastings05_hierarchical_projection_projection_expectation}) alone, the recursion would have $\sqrt{b_{m-1}}$ in
place of $b_{m-1}$.

The spectral step proves more than~(\ref{eq:hastings05_hierarchical_projection_projection_expectation}). Write
$N=M-c$. On the range of $1-P^1$ every eigenvalue of $N$ has absolute value
greater than $x$, so
\begin{align}\label{eq:hastings05_hierarchical_projection_projection_residual}
    \lVert\Psi-P^1\Psi\rVert\le\lVert N\Psi\rVert/x\le\epsilon/x .
\end{align}
The formalization measures the deficiency of a level in norm,
$b=\lVert\Psi-P\Psi\rVert$. This implies~(\ref{eq:hastings05_hierarchical_projection_projection_expectation}): with
$d=\lVert\Psi-P\Psi\rVert\le1$,
\begin{align}\label{eq:hastings05_hierarchical_projection_expectation_norm}
    \langle\Psi,P\Psi\rangle=1-d^2\ge1-d .
\end{align}
With $b$ in norm,~(\ref{eq:hastings05_hierarchical_projection_normalized_projection_distance}) holds as printed. For commuting
orthogonal projections $P,P'$ the vectors $\Psi-P\Psi$ and $P(\Psi-P'\Psi)$
are orthogonal, so
$\lVert\Psi-PP'\Psi\rVert^2\le\lVert\Psi-P\Psi\rVert^2+\lVert\Psi-P'\Psi\rVert^2\le2b^2$.
For an orthogonal projection $Q$ and a unit vector $\Psi$ with
$t=\lVert Q\Psi\rVert>0$,
\begin{align}\label{eq:hastings05_hierarchical_projection_normalization_sq}
    \Bigl\lVert\frac{Q\Psi}{t}-\Psi\Bigr\rVert^2=2-2t\le2(1-t^2)
    =2\lVert\Psi-Q\Psi\rVert^2 .
\end{align}
With $Q=PP'$, (\ref{eq:hastings05_hierarchical_projection_normalization_sq}) and
$\lVert\Psi-PP'\Psi\rVert^2\le2b^2$ give
$\lVert\Psi^1_a-\Psi_a\rVert\le2b$, so
the residual at the second level and the printed $b_2$ follow as in the
source. If the projected vector vanishes, normalization is interpreted as
zero and the corresponding inequality follows directly.

From the third level on, the window stays centred at the expectation in
$\Psi_a$. Thus the residual of a later normalized vector must be compared
with $\Psi_a$. If consecutive normalized vectors differ by at most $2b_j$,
the triangle inequality gives
$\lVert\Psi^{m-1}_a-\Psi_a\rVert\le2\sum_{j<m}b_j$. Under this assumption,
an accumulated scalar majorant has the recurrence
\begin{align}\label{eq:hastings05_hierarchical_projection_accumulated_recurrence}
    \beta_0=0,\qquad
    \beta_m=\Bigl(\epsilon+2\sum_{1\le j<m}\beta_j\Bigr)^{1-1/h}.
\end{align}
For $0<\epsilon\le1$ and $h\ge1$, the coefficient estimate
$\beta_m\le c_m\epsilon^{(1-1/h)^m}$ has
$c_1=1$ and $c_m=1+2\sum_{1\le j<m}c_j=3^{m-1}$.
This remains polynomial in $N=2^h$. It does not prove the printed
recurrence with only the preceding deficiency, nor rule out a sharper
argument giving that recurrence. The Lean
recurrence theorems establish inequalities for the printed scalar sequence;
they do not identify it with deficiencies beyond the second level. The
complete hierarchy is instead controlled by measuring every deficiency
against $\Psi_a$; see
Section~\ref{sec:hastings05_hierarchical_projection_all_levels}.

\newpage
\section{The restricted spectral window}

At the second level the source takes eigenvectors of $QMQ$ inside
$\operatorname{ran}Q$, where $Q=P^1_kP^1_{k+2}$
\cite[lines 673--685]{Hastings2006Solving}.
Let $W=P_x(QMQ,c)$ denote the spectral window on the whole ambient space.
Since $QMQ$ vanishes on $\ker Q$, the projection $W$ includes $\ker Q$
whenever $|c|\le x$. Thus the later-level projection is $P=QW$, not $W$.
For example, with $Q=\operatorname{diag}(1,0)$, $M=0$, $c=0$ and $x=1$,
one has $W=1$ whereas $P=Q$.

The matrices $Q$ and $QMQ$ commute, so functional calculus gives $QW=WQ$.
Consequently $P$ is an orthogonal projection with $PQ=P$ and annihilates
$\ker Q$. On a vector $\phi$ with $Q\phi=\phi$, it agrees with $W$, and
\begin{align}\label{eq:hastings05_hierarchical_projection_restricted_error}
    \lVert\phi-P\phi\rVert
     & \le\frac{\lVert(QMQ-c)\phi\rVert}{x}
    \le\frac{\lVert(M-c)\phi\rVert}{x}.
\end{align}
Here $x>0$, and the second inequality in
(\ref{eq:hastings05_hierarchical_projection_restricted_error}) follows from
$(QMQ-c)\phi=Q(M-c)\phi$. This implements the subspace specified by the
source without imposing $|c|>x$. It establishes the required nesting of a
single level, but does not establish the full hierarchy or its recurrence.

\section{The residual and the recursion}

The centres in the source's second-level calculation differ: the spectral
window is centred at $c=\langle\Psi_a,M\Psi_a\rangle$ (line 681), whereas
the displayed residual uses $c'=\langle\Psi^1_a,M\Psi^1_a\rangle$
(lines 692--694). A residual bound about $c'$ alone does not give the same
bound about $c$: the triangle inequality adds $|c'-c|$ for a unit vector.
The formalized perturbation estimate instead keeps $c$ fixed throughout,
as required by that spectral window. Its coefficient is $\lVert M-c\rVert$;
identifying it with the printed $J l_{\mathrm{proj}}$ requires the
corresponding operator-norm bound.

For any vectors $\phi,\psi$,
$\lVert N\phi\rVert\le\lVert N\psi\rVert+\lVert N\rVert\,\lVert\phi-\psi\rVert$,
and $\lVert N\rVert\le\lVert M\rVert+|c|\le2\lVert M\rVert$ since
$|c|\le\lVert M\rVert$. The source's term $2b_1J\,l_{\mathrm{proj}}$ reads
$\lVert N\rVert\le J\,l_{\mathrm{proj}}$. The formalization keeps
$\lVert N\rVert$ in the bound and states the recursion with a general
coefficient $a\ge0$:
\begin{align}\label{eq:hastings05_hierarchical_projection_scalar_recurrence}
    b_0=0,\qquad b_{m+1}=(\epsilon+ab_m)^s,\qquad
    b_m\le\Bigl(\sum_{j<m}a^j\Bigr)\epsilon^{s^m}
\end{align}
for $0<\epsilon\le1$ and $0\le s\le1$. The induction uses
$\epsilon\le\epsilon^{s^m}$ and $C^s\le C$ for $C\ge1$. For the source's $a=2$
and $s=1-1/h$,
\begin{align}\label{eq:hastings05_hierarchical_projection_scalar_recurrence_bound}
    b_h\le(2^h-1)\,\epsilon^{(1-1/h)^h},
\end{align}
which for $N=2^h$ and $\epsilon=e^{-l_{\mathrm{proj}}/l_{\mathrm{fac}}}$ is
below the source's $N\exp[-(l_{\mathrm{proj}}/l_{\mathrm{fac}})(1-1/h)^h]$
(lines 705--707). No constant is lost. The source continues with
$\sim N\exp(-l_{\mathrm{proj}}/(e\,l_{\mathrm{fac}}))$. Since $(1-1/h)^h$
increases to $1/e$, this is the limiting order of magnitude and not an upper
bound. The formalization proves $(1-1/h)^h\ge1/4$ for $h\ge2$, so
$b_h\le(2^h-1)\epsilon^{1/4}$.

\section{The energy estimate and higher-dimensional grouping}

The energy bound eq.~(hchy) and the count of at most $2n$ final states
(lines 709--730) are proved for the corrected periodic construction in the
\href{hastings05_hierarchical_energy.pdf}{hierarchical-energy note}.
That construction now starts from the original physical interactions and
uses an explicit integer quadratic logarithmic radius at every physical
chain length. The coordinate-slice option in higher dimensions
(lines 745--756) is constructed on full finite rectangular boxes in the
\href{hastings05_coordinate_slices.pdf}{coordinate-slices note}; its
scope and polynomial logarithmic range are stated there.

\section{All levels against the ground state}
\label{sec:hastings05_hierarchical_projection_all_levels}

The source projects at level $m$ inside the
range of the two projections of level $m-1$ (lines 671--676). For orthogonal
projections $P\le Q$ and a unit vector $\Psi$ with
$\Psi'=Q\Psi/\lVert Q\Psi\rVert$, the formalization proves
\begin{align}\label{eq:hastings05_hierarchical_projection_nested_deficiency}
    \lVert\Psi-P\Psi\rVert^2\le\lVert\Psi-Q\Psi\rVert^2
    +\lVert\Psi'-P\Psi'\rVert^2,
\end{align}
which combines the deficiencies of consecutive levels. Along such a chain the
deficiency of $\Psi_a$ under a projection of level $m$ collects the squared
deficiencies of the earlier levels in addition to $b_m^2$.

The formalization therefore measures every deficiency against $\Psi_a$
itself. Let $P_k^m$ be the projection of level $m$ on block $k$, with
$P_k^0=1$, and let $P_k^{m+1}$ be the window of the pair operator $M$ joining
the two halves of block $k$, centred at $c=\langle\Psi_a,M\Psi_a\rangle$ with
radius $x_m$, inside the range of $Q=P_{2k}^mP_{2k+1}^m$. The two halves act
on disjoint supersites, so their projections commute. Suppose
$\lVert(M-c)\Psi_a\rVert\le\epsilon$ and $\lVert M-c\rVert\le K$ for every
pair operator, and let each half move $\Psi_a$ by at most $D_m$. Then
$\lVert\Psi_a-Q\Psi_a\rVert^2\le2D_m^2$, and the normalized image
$\Psi'=Q\Psi_a/\lVert Q\Psi_a\rVert$ lies in the range of $Q$ at distance at
most $2D_m$ from $\Psi_a$. By the restricted window bound,
$\lVert\Psi'-P_k^{m+1}\Psi'\rVert\le(\epsilon+2KD_m)/x_m$, and
(\ref{eq:hastings05_hierarchical_projection_nested_deficiency}) gives
\begin{align}\label{eq:hastings05_hierarchical_projection_ground_recurrence}
    D_0=0,\qquad
    D_{m+1}^2\le2D_m^2+\Bigl(\frac{\epsilon+2KD_m}{x_m}\Bigr)^2 .
\end{align}
With $K>0$, $e>0$, $\epsilon=Ke$ and the source's radius
$x_m=K(e+2\beta_m)^{1-s}$ (lines 676--681), where $\beta_0=0$ and
$\beta_{m+1}^2=2\beta_m^2+(e+2\beta_m)^{2s}$, induction gives
$D_m\le\beta_m$. For $0<e\le1$ and $0\le s\le1$,
\begin{align}\label{eq:hastings05_hierarchical_projection_ground_bound}
    \beta_m\le(3^m-1)\,e^{s^m},
\end{align}
since $C_{m+1}^2=2C_m^2+(1+2C_m)^2\le(3C_m+2)^2$ for the coefficients. This
replaces the printed recurrence and its coefficient $2^h-1$ by
(\ref{eq:hastings05_hierarchical_projection_ground_recurrence}) and $3^h-1$,
which is still polynomial in $N=2^h$. With $s=1-1/h$ and $h\ge2$ the final
deficiency is at most $(3^h-1)e^{1/4}$.

\section{Uniform radii and an explicit projection length}

Put $r=l_{\mathrm{proj}}/l_{\mathrm{fac}}$, $e=e^{-r}$ and
$s=1-1/h$, with $h\ge2$. The nonnegative sequence $\beta_m$ is
increasing, since its squared recurrence dominates $2\beta_m^2$.
Consequently, for $m\le h$,
\[
 \beta_m\le(3^h-1)e^{-r/4}.
\]
The adaptive radius therefore satisfies
\[
 x_m=K(e+2\beta_m)^{1/h}
 \le K\bigl[(2\cdot3^h-1)e^{-r/4}\bigr]^{1/h}
 \le9K e^{-r/(4h)}.
\]
Every radius is strictly positive for $K>0$. The deficiency after the
$h$ tree levels is less than one if $r>4h\log3$.

To control the periodic closing link, add its spectral window inside
the last retained space, as explained in
\href{hastings05_hierarchical_energy.pdf}{the hierarchical-energy note}.
The same deficiency majorant bounds this extra step. Since
$(1-1/h)^{h+1}\ge(1/4)(1/2)=1/8$, its deficiency is at most
\[
 (3^{h+1}-1)e^{-r/8},
\]
and hence is less than one if $r>8(h+1)\log3$.
These statements concern the explicitly corrected norm deficiency,
not the printed expectation-deficiency recurrence.\footnote{Declarations:
\leanid{TNLean.MPS.AreaLaw.LocalSolution.hierarchyWindowRadius_le},
\leanid{TNLean.MPS.AreaLaw.LocalSolution.hierarchyDeficiency_window_lt_one},
\leanid{TNLean.MPS.AreaLaw.LocalSolution.hierarchyDeficiency_window_succ_le}, and
\leanid{TNLean.MPS.AreaLaw.LocalSolution.hierarchyDeficiency_window_succ_lt_one}.}

For the scalar estimate, let $J_{\mathrm{eff}}>0$ and suppose
$K\le J_{\mathrm{eff}}r$. In the actual one-dimensional periodic construction
this inequality is derived, with $K=8\ell J$, $r=\eta\ell$ and
$J_{\mathrm{eff}}=8J/\eta$, as detailed below. Let
$\Delta>0$ be the physical spectral gap. Define
\[
 A=1+\max\{0,\log(J_{\mathrm{eff}}/\Delta)\},\qquad
 r_*=512h^2 A.
\]
For every $r\ge r_*$, all $2^h-1$ tree windows and the closing window
obey
\begin{align*}
 \sum_{m=0}^{h-1}2^{h-1-m}x_m+x_h
   &\le9\cdot2^hK e^{-r/(4h)}\le\Delta/4,\\
 2^hJ_{\mathrm{eff}}r e^{-r}&\le\Delta/8.
\end{align*}
The length $r_*$ also exceeds the strict closing-window threshold above.
In particular the retained space is nonzero, and each radius is at most
$\Delta/(4\cdot2^h)$. This proves an explicit version of the source's
quadratic dependence on $h=\log_2N$ in lines~709--720; the fixed local
constants are included transparently in $J_{\mathrm{eff}}$.

For a chain with an arbitrary finite number $N$ of supersites, set
$h=\max\{2,\lceil\log_2 N\rceil\}$, with the upper integer logarithm
read as zero for $N\le1$. Extend the finite edge family by zero up to
$2^h$ edges. This changes neither its physical sum nor the local
residual and centred-norm estimates. The chosen depth satisfies
$N\le2^h$ and $h\le\lceil\log_2N\rceil+2$, so the explicit projection
length is bounded by
\[
 512(\lceil\log_2N\rceil+2)^2
 \bigl(1+\max\{0,\log(J_{\mathrm{eff}}/\Delta)\}\bigr).
\]
No assumption that the physical chain length is a power of two is
needed.\footnote{Declarations:
\leanid{TNLean.MPS.AreaLaw.LocalSolution.finiteHierarchyDepth},
\leanid{TNLean.MPS.AreaLaw.LocalSolution.le_two_pow_finiteHierarchyDepth},
\leanid{TNLean.MPS.AreaLaw.LocalSolution.finiteHierarchyDepth_le_clog_add_two}, and
\leanid{TNLean.MPS.AreaLaw.LocalSolution.hierarchyProjectionLength_finiteHierarchyDepth_le}.}

For completeness, the scalar estimate uses $t e^{-t}\le1$ to absorb
the linear factor in $K$:
\[
 r e^{-r/(4h)}\le8h e^{-r/(8h)}.
\]
Since $2^h\le e^h$, $h\le e^h$, $72\le e^{36h}$ for $h\ge2$, and
$J_{\mathrm{eff}}\le\Delta e^A$, the energy envelope is at most
\[
 \Delta\exp\bigl(A+38h-r/(8h)\bigr)
 \le\Delta e^{-2}\le\Delta/4.
\]
The exponential truncation estimate follows from the same scalar result
with coupling bound $2J_{\mathrm{eff}}r/9$, which is at most
$J_{\mathrm{eff}}r$, and $e^{-r}\le e^{-r/(4h)}$.
Thus no final energy-smallness assumption has been inserted.\footnote{Declarations:
\leanid{TNLean.MPS.AreaLaw.LocalSolution.hierarchyProjectionLength},
\leanid{TNLean.MPS.AreaLaw.LocalSolution.hierarchyWindowEnergy_le_quarter_gap},
\leanid{TNLean.MPS.AreaLaw.LocalSolution.hierarchyDeficiency_window_succ_lt_one_of_length},
\leanid{TNLean.MPS.AreaLaw.LocalSolution.hierarchySum_windowRadius_add_le_quarter_gap}, and
\leanid{TNLean.MPS.AreaLaw.LocalSolution.hierarchyTruncationError_le_eighth_gap}, in
\doclink{TNLean/MPS/AreaLaw/LocalSolution/HierarchyRadius}, and
\leanid{TNLean.MPS.AreaLaw.LocalSolution.hierarchyTruncationError_mul_length_le_eighth_gap}.}

\subsection{The actual one-dimensional periodic bounds}

For a physical cycle of $M\ge1$ sites, take
$N=\max\{1,\lfloor M/(2\ell)\rfloor\}$ contiguous blocks, putting the
remainder in the last block. Every cyclic radius-$\ell$ term belongs to
an adjacent pair, and at most $4\ell$ terms are assigned to a pair.
Consequently local norm $J$ and residual $Je^{-\eta\ell}$ give the actual
centred pair bound $K=8\ell J$ and residual at most $Ke^{-\eta\ell}$.
Writing $r=\eta\ell$, the effective coupling is $J_{\rm eff}=8J/\eta$,
independent of physical chain length. Equal blocks improve the constant
$8$ to $4$; they are not required.

The integer choice
\begin{align*}
 \ell=\left\lceil\frac{512}{\eta}(\lceil\log_2M\rceil+2)^2
      \bigl(1+\max\{0,\log((8J/\eta)/\Delta)\}\bigr)\right\rceil
\end{align*}
is positive and satisfies the hierarchy length inequality for the block
count determined by this very $\ell$. It is less than one plus the
quantity inside the ceiling. This proves the quadratic logarithmic
physical range without a divisibility condition or a supplied pair-norm
hypothesis. For a physical interaction range $R$, take the maximum of
this integer with $\lceil R\rceil$; the explicit upper bound increases
by at most $\lceil R\rceil$, and the sufficient hierarchy length still
holds. The actual cyclic metric has ball count at most $2\lceil R\rceil+1$;
its localization theorem constructs the required truncated operators,
including the exact zero-range case. The physical approximation uses
$\sum_i\widetilde H_i=H$ and the small differences
$\widetilde H_i-M_i$, where $\widetilde H_i$ denotes the full Gaussian
filter, as explained in the
\href{hastings05_hierarchical_energy.pdf}{hierarchical-energy note}.
The local residual condition (low) is retained, and the closing-window
correction preserves every lower retained space.\footnote{Declarations:
\leanid{TNLean.MPS.AreaLaw.LocalSolution.centeredNorm_roundedPeriodicOperator_le},
\leanid{TNLean.MPS.AreaLaw.LocalSolution.residual_roundedPeriodicOperator_le},
\leanid{TNLean.MPS.AreaLaw.LocalSolution.roundedPeriodicProjectionRadius},
\leanid{TNLean.MPS.AreaLaw.LocalSolution.hierarchyProjectionLength_le_chosen_roundedPeriodicRadius},
\leanid{TNLean.MPS.AreaLaw.LocalSolution.exists_unit_mem_roundedPeriodicProjection_of_chosen_radius}, and
\leanid{TNLean.MPS.AreaLaw.LocalSolution.card_le_two_mul_re_trace_of_roundedPeriodicProjection_of_chosen_radius},
\leanid{TNLean.MPS.AreaLaw.LocalSolution.physicalPeriodicProjectionRadius_lt}, and
\leanid{TNLean.MPS.AreaLaw.LocalSolution.exists_periodicPhysicalHierarchy}.}

\subsection{An effective coupling depending on the system size}

A coarse pair-operator norm bound may make the effective coupling
proportional to the total number of physical terms. This still preserves
polynomial logarithmic growth of the projection length. Precisely, for
$N\ge1$, $C,J,\Delta>0$ and $J_{\mathrm{eff}}=NCJ$, put
$L=\lceil\log_2N\rceil$ and $A=1+\max\{0,\log(CJ/\Delta)\}$. Then
\[
 1+\max\{0,\log(J_{\mathrm{eff}}/\Delta)\}
 \le\log N+A\le(L+2)A.
\]
With the chosen depth $h=\max\{2,L\}$ this gives the proved bound
\[
 r_*\le512(L+2)^3 A.
\]
Thus a linear-in-size effective coupling gives a cubic logarithmic
scalar length at fixed local coupling and gap. This is relevant to the
higher-dimensional extension in lines~745--756; applying it to a
particular grouping requires its physical support and residual estimates.
These are constructed for the coordinate-slice option on full finite
rectangular boxes in every dimension; see
\href{hastings05_coordinate_slices.pdf}{the coordinate-slices note}.\footnote{Declaration:
\leanid{TNLean.MPS.AreaLaw.LocalSolution.hierarchyProjectionLength_finiteHierarchyDepth_mul_le}.}

\section{Formalization}

The spectral residual bound (\ref{eq:hastings05_hierarchical_projection_projection_residual}) and its expectation consequence
(\ref{eq:hastings05_hierarchical_projection_projection_expectation}) are formalized, as is the normalization bound
(\ref{eq:hastings05_hierarchical_projection_normalized_projection_distance}) with deficiency measured in norm.\footnote{Declarations:
    \leanid{TNLean.MPS.AreaLaw.LocalSolution.norm_sub_windowProjection_le},
    \leanid{TNLean.MPS.AreaLaw.LocalSolution.one_sub_div_le_re_inner_windowProjection}, and
    \leanid{TNLean.MPS.AreaLaw.LocalSolution.norm_inv_norm_smul_sub_le_two_mul}.}

The restricted spectral window $P=QP_x(QMQ,c)$ satisfies $PQ=P$.
Its residual estimate combines with the distance to a reference vector
to give the second-level bound.\footnote{Declarations:
    \leanid{TNLean.MPS.AreaLaw.LocalSolution.compressedWindowProjection},
    \leanid{TNLean.MPS.AreaLaw.LocalSolution.compressedWindowProjection_mul}, and
    \leanid{TNLean.MPS.AreaLaw.LocalSolution.norm_sub_compressedWindowProjection_le_of_norm_sub_le}.}

The bounds for the scalar recurrence in (\ref{eq:hastings05_hierarchical_projection_scalar_recurrence}) and
(\ref{eq:hastings05_hierarchical_projection_scalar_recurrence_bound}) are formalized independently of the complete
hierarchy.\footnote{Declarations:
    \leanid{TNLean.MPS.AreaLaw.LocalSolution.levelDeficiency_le} and
    \leanid{TNLean.MPS.AreaLaw.LocalSolution.levelDeficiency_self_le}.}

The bound over all levels,
(\ref{eq:hastings05_hierarchical_projection_ground_recurrence}) and
(\ref{eq:hastings05_hierarchical_projection_ground_bound}), is
formalized.\footnote{Declarations:
    \leanid{TNLean.MPS.AreaLaw.LocalSolution.norm_sub_hierarchy_le},
    \leanid{TNLean.MPS.AreaLaw.LocalSolution.hierarchyDeficiency_le_hierarchyMajorant}, and
    \leanid{TNLean.MPS.AreaLaw.LocalSolution.hierarchyMajorant_le}.}
The commutation of the two halves of each block, which the bound over all
levels takes as a hypothesis, is derived from the supersite grouping of lines
609--628: a spectral window keeps the support of its operator, so every
projection of level $m$ acts on its block of $2^m$ supersites, and the two
halves of a block act on disjoint sites. The residual of a pair operator is
at most the sum of the residuals of its summands $M_i$, with centre the sum of
their expectations.\footnote{Declarations:
    \leanid{TNLean.MPS.AreaLaw.LocalSolution.norm_sub_supersiteHierarchy_le},
    \leanid{TNLean.MPS.AreaLaw.LocalSolution.mem_supportedOperators_hierarchy}, and
    \leanid{TNLean.MPS.AreaLaw.LocalSolution.norm_toEuclideanLin_sum_sub_le}.}

\bibliographystyle{plain}
\bibliography{references}

\end{document}
